% !PRINTED 232
\weberpage{232}
%<<CONT
and denote by $h$, $k$ two constants, then
\[
\Phi(y) = h\,\Phi(x), \quad \Psi(y) = k\,\Psi(x),
\]
provided that $\begin{pmatrix} \alpha, & \beta \\ \gamma, & \delta \end{pmatrix}$ is one of the icosahedral substitutions. But now, since $y$ as well as $x$ is an independent variable, by (7) the identity
\[
\Phi(x) + \Psi(x) = h\,\Phi(x) + k\,\Psi(x) = 1728,
\]
holds, and since $\Phi$ and $\Psi$ do not stand in a constant ratio, we must have $h = k = 1$; that is, the quotients $\Phi$ and $\Psi$ remain absolutely unchanged under the icosahedral substitutions.

From this it follows further \gesperrt{that $f$, $H$, $T$ remain absolutely unchanged under the homogeneous icosahedral substitutions with determinant $1$.}

For if $f$ passes under such a substitution into $cf$, then $H$ and $T$ pass into $c^2H$ and $c^3T$ [by (2) and (4)], and the quotients $\Phi$ and $\Psi$ into $c\Phi$ and $c\Psi$. Since, on the other hand, $\Phi$ and $\Psi$ remain unchanged, $c = 1$.

If we set $\Psi(x) = z$, then $\Phi(x) = 1728 - z$, and from (8) we obtain the two equations:
\weq{9}{H^3 - z f^5 = 0,}
\weq{10}{T^2 - (1728 - z)\,f^5 = 0,}
of which, on account of the identity (6), the one follows from the other, so that they are really only two different forms of one and the same equation. If we now regard $z$ in it as given, we have an equation of degree 60 for $x$, which is called the \gesperrt{icosahedral equation}.

We shall return to the properties of this equation and to its relation to the general equation of degree 5 in a later section.

\begin{center}
\S.\,61.\\[2pt]
\gesperrt{The Invariants of the Icosahedron.}
\end{center}

By the fundamental forms of the polyhedral groups with which we have so far become acquainted, the region of the invariants of the groups in question is, as we can now prove, exhausted. We
%CONT>>
